
    \section{Real-World System Implementation and Experiments}
    We built a prototype real-world robotic system on Linux to evaluate the proposed scheduler.
    The system runs two types of periodic tasks: robotic application tasks and a scheduler task.

    \subsection{Scheduler Implementation Design}
    In each call, the scheduler task (i) collects recently recorded execution-time data from each robotic task, 
	(ii) predicts execution-time distributions for the next period, 
	(iii) optimizes priority assignments and QoS budgets, and (iv) updates the tasks' priority assignments and configurations.
    All computations are performed in user space rather than kernel space to simplify implementation; although this introduces higher overhead, it remains lightweight and practical because of our lightweight optimization algorithm design.
    

    \subsection{Computation tasks}
    \label{section_computation_tasks}
    All tasks are modeled as periodic tasks.
    A new job of a task $\tau_i$ is released for execution if the previous released job instance has completed execution and the new job's release time has passed.
    Jobs that miss their deadlines are not discarded.
    Following the approach in Sifat~\etal~\cite{Ash23RAL}, our system includes the following robotic tasks:

    \textit{Simultaneous Localization and Mapping (SLAM):}
    We based our implementation on ORB-SLAM2~\cite{murORB2}, a well-known real-time visual SLAM system.
    % The original ORB-SLAM2 implementation uses multiple threads for sub-tasks (\eg tracking, loop closure, and visualization), some of which involve infinite loops.
    To align with our real-time systems model, we modified the SLAM implementation without influencing its functionality: 
	we replaced infinite loops with periodic tasks, changed the original multi-thread implementation into a single main thread as embedded systems typically have limited resources, 
	and disabled real-time map visualization.
    These modifications preserve core functionality while improve execution-time predictability.
    If multi-threaded implementation is required, techniques such as gang task scheduling and response time analysis models~\cite{Sun2024StrictPF} can be used to analyze response time distributions.
    % --- multi-core (gang) task scheduling: retained as comments for future reference ---
    % \subsection{Multi-core Tasks Scheduling}
    % \label{section: gang}
    % In Section~\ref{section: system_model_tasks}, we assumed that each computation task only occupies one single CPU core during its execution.
    % Such kind of task model includes both single-thread tasks and some multi-thread tasks that do not need to run more than one threads in parallel.
    % These tasks are very similar to single-thread tasks from the scheduling purposes.
    %
    % Other types of multi-thread tasks may include those that need to run more than one threads in parallel. These tasks can be modelled as gang tasks~\cite{Goossens2010GangFS}.
    % The proposed optimization framework in this paper can work with gang tasks if their response time distribution analysis is available. However, such analysis could be complicated and computationally expensive.
    %
    % Therefore, to incorporate the gang tasks into our system, it is easier to treat the gang tasks as single-thread tasks from the scheduling perspective. This can be done by ``merging'' the multiple computation cores required by the gang tasks into one computation core during scheduling.
    % These ``merged'' cores will not be allocated to more than one task each time, and therefore behave similarly to a single core.
    % This method works well for GPU tasks that run on embedded systems, one of the major applications of gang tasks. This is because that GPU tasks typically cannot run in parallel due to the high context switch overhead in low-end GPUs~\cite{Olmedo2020DissectingTC} in embedded systems.

    \textit{Traveling Salesperson Problem (TSP):}
    We adopted the anytime algorithm implementation from Huang~\etal~\cite{Huang23ICRA_TSP}.
    The running time limit for TSP is treated as a configurable optimization variable in our system, where longer time limits yield shorter output paths.

    \textit{Path Planning:}
    Our implementation of the rapidly-exploring random tree (RRT) algorithm~\cite{LaValle1998RapidlyexploringRT} is based on~\cite{rrt_github}.
    The RRT's execution time varies based on input map configurations and the algorithm's inherent randomness.

    \textit{Model Predictive Controller (MPC):}
    We implemented a constrained linear-quadratic MPC using the OSQP solver~\cite{stellato18OSQP}.
    MPC exhibits relatively stable execution times.

	To enhance the system's predictability, input data is pre-loaded into memory whenever possible without violating the experiments' authenticity, reducing potential I/O contention and variations in execution-time distribution.
	
	\subsection{Real-World Experiment Parameters}
        
\todo{Update the repository name later.}

    We tested the scheduler on an NVIDIA Jetson AGX Xavier (Ubuntu 18), restricting all computation tasks to two of its eight cores to simulate a resource-constrained platform with heavier computation loads.
    % This setting also artificially simulates a system with more computation loads that do not have much flexibility in scheduling, such as computer vision tasks that require delicate collaboration between CPU and GPU.
    A third core was dedicated to the scheduler for precise overhead measurement, though this is unnecessary in practice given the scheduler's light computation.
    All results are reproducible with our released code\footnote{\githubLink{}}.

    The runtime configuration parameters, summarized in Table~\ref{table_exp_config}, are mainly adopted from Sifat~\etal~\cite{Ash23RAL} with adjustments based on the tasks' semantics.
    Safety-critical tasks (MPC, SLAM) receive higher deadline-miss thresholds and weights than TSP and RRT, with SLAM weighted highest as accurate localization is safety-critical.
    % Such a setting also simulates a realistic task set with a mix of critical and non-critical tasks.
    Only TSP uses an anytime algorithm, so it is the only task with a non-constant performance metric;
    Periods are higher than each task's worst-case execution time and are rounded to convenient values, with some (e.g., SLAM) larger than common choices due to the high execution time in current implementation, 
	and engineering highly efficient implementations is not this paper's focus.
    % Original Period detail: periods were rounded to convenient, "nice" values (e.g., integers), as is common in real-world systems. Since developing highly efficient application algorithms and implementation is not the focus of this paper, some applications, such as SLAM, may have a larger period compared with real-world applications.
    The MPC task is considered as the safety-critical task in experiments.
    
	\begin{table}[htbp]
    \centering
    \caption{Experiment Configurations}
    \label{table_exp_config}
\begin{tabular}{@{}llllll@{}}
\toprule
     & \begin{tabular}[c]{@{}l@{}}Miss deadline \\ threshold\end{tabular} & Weight & \begin{tabular}[c]{@{}l@{}}Performance\\ Function\end{tabular} & Period (ms) & \begin{tabular}[c]{@{}l@{}}CPU \\ ID\end{tabular} \\ \midrule
SLAM & 0.9                                                                & 2      & Constant                                                       & 2000        & 0                                                 \\
TSP  & 0.5                                                                & 1      & Varied                                                         & 1500        & 0                                                 \\
RRT  & 0.5                                                                & 1      & Constant                                                       & 4000        & 1                                                 \\
MPC  & 0.99                                                               & 1      & Constant                                                       & 10          & 1                                                 \\ \bottomrule
\end{tabular}
\end{table}
	
	\subsection{Scheduling algorithms in experiments}
    \label{section: exp_methods}
    We evaluate five algorithms, used in both the real-world experiments here
    and the simulation experiments (Section~\ref{section: simulation}):
	\begin{itemize}
		\item \textbf{CFS}: the default Linux Completely Fair Scheduler.
        It is not a real-time scheduler and provides no performance guarantees.
		\item \textbf{DM\_FAST} and \textbf{DM\_SLOW}: Deadline-Monotonic
        fixed-priority schedulers (Definition~\ref{def_dm_fast}) that fix every
        QoS budget: shortest for DM\_FAST, longest for DM\_SLOW.
		\item \textbf{BF}: the brute-force optimizer
        (Sections~\ref{section_bf_pa}, \ref{section_config_opt}) as a baseline.
		\item \textbf{INCR}: the proposed incremental optimizer
        (Sections~\ref{section_increment_pa}, \ref{section_config_opt}),
        which jointly optimizes priority assignments and QoS budgets.
	\end{itemize}
	
    To reduce the influence of random factors and improve reproducibility,
    all results are based on $50$ runs under identical conditions and
    reported as averages or box plots.
	\begin{figure*}[ht!]
		\centering 
		
		\begin{subfigure}[t]{0.32\textwidth} 
			\centering
			\includegraphics[width=\textwidth]{pictures/box_plot_ET_data_SLAM_optimizerIncremental.pdf}
			\caption{SLAM, incremental optimizer}
			\label{fig_box_et_SLAM}
		\end{subfigure}
		\begin{subfigure}[t]{0.32\textwidth} 
			\centering
			\includegraphics[width=\textwidth]{pictures/box_plot_ET_data_TSP_optimizerIncremental.pdf}
			\caption{TSP, incremental scheduler}
			\label{fig_box_et_TSP}
		\end{subfigure}
		\begin{subfigure}[t]{0.32\textwidth} 
			\centering
			\includegraphics[width=\textwidth]{pictures/box_plot_ET_data_RRT_optimizerIncremental.pdf}
			\caption{RRT}
			\label{fig_box_et_RRT}
		\end{subfigure}
		
		\begin{subfigure}[t]{0.32\textwidth} 
			\centering
			\includegraphics[width=\textwidth]{pictures/box_plot_ET_data_MPC_optimizerIncremental.pdf}
			\caption{MPC}
			\label{fig_box_et_MPC}
		\end{subfigure}
		\begin{subfigure}[t]{0.32\textwidth} 
			\centering
			\includegraphics[width=\textwidth]{pictures/box_plot_ET_data_SCHEDULER_PUER_OPT_optimizerBF.pdf}
			\caption{Brute-Force Scheduler}
			\label{fig_box_et_bf}
		\end{subfigure}
		\begin{subfigure}[t]{0.32\textwidth} 
			\centering
			\includegraphics[width=\textwidth]{pictures/box_plot_ET_data_SCHEDULER_PUER_OPT_optimizerIncremental.pdf}
			\caption{Incremental scheduler}
			\label{fig_box_et_incremental}
		\end{subfigure}
		\caption{Execution time variation at different time during experiments} 
		\label{fig_et_all}
	\end{figure*}

    
	\begin{figure*}[ht!]
		\centering 
		\begin{subfigure}[t]{0.32\textwidth} 
			\centering
			\includegraphics[width=\textwidth]{pictures/box_plot_of_all_data_CFS.pdf}
			\caption{CFS}
			\label{fig_cfs_perf}
		\end{subfigure}
		\begin{subfigure}[t]{0.32\textwidth}
			\centering
			\includegraphics[width=\textwidth]{pictures/box_plot_of_all_data_RM_Fast.pdf}
			\caption{DM\_Fast}
			\label{fig_dm_fast}
		\end{subfigure}
		\begin{subfigure}[t]{0.32\textwidth} 
			\centering
			\includegraphics[width=\textwidth]{pictures/box_plot_of_all_data_RM_Slow.pdf}
			\caption{DM\_Slow}
			\label{fig_dm_slow}
		\end{subfigure}
		
		\begin{subfigure}[c]{0.32\textwidth}
			\centering
			\includegraphics[width=\textwidth]{pictures/box_plot_of_all_data_optimizerBF.pdf}
			\caption{Brute-Force Scheduler}
			\label{fig_opt_bf}
		\end{subfigure}
		\begin{subfigure}[c]{0.32\textwidth}
			\centering
			\includegraphics[width=\textwidth]{pictures/box_plot_of_all_data_optimizerIncremental.pdf}
			\caption{Incremental Scheduler}
			\label{fig_opt_incre}
		\end{subfigure}
		% \begin{subfigure}[c]{0.32\textwidth}
		% 	\centering
  %           \includegraphics[width=\textwidth]{pictures/average_all.pdf}
		% 	\caption{Avergage of all methods}
		% 	\label{fig_avg_all}
		% \end{subfigure}
		
		\caption{The SP metric of different scheduling methods} 
		\label{figs_exp_results_all_cps}
	\end{figure*}

	\subsection{Execution-time dynamics during the experiments}
    Fig.~\ref{fig_et_all} reports per-task execution-time box plots over the
    experiments.
    Execution times are measured as CPU time, excluding preemption; for the
    scheduler, where CPU time is hard to measure precisely, we use end-to-end
    time as a conservative upper bound.
    (In Fig.~\ref{fig_box_et_TSP}, some TSP instances are not run, yielding
    zero-execution-time points.)

    Together these plots show how the workload evolves and how each scheduler
    responds: MPC and RRT execution times are relatively stable, with RRT showing
    slightly higher variance, while the notable dynamics are in SLAM and TSP.

    SLAM's execution time gradually rises and then falls over the course of a run.
    The rise reflects growing computational demand at different time: loop closure adds work, and
    optimizing over an accumulating set of frames costs more, until the system
    reaches a state where demand drops again.

    TSP's execution time is set by its runtime time limit.
    Under CFS, DM\_FAST, and DM\_SLOW this limit is fixed, so TSP's distribution
    stays stable; these schedulers never change it.
    Under the brute-force and incremental optimizers, TSP's limit varies over time
    as the scheduler reallocates resources in response to other tasks.
    Around $650$~s, for instance, as SLAM's demand grows, the optimizer trims TSP's
    budget so that the higher-priority SLAM meets its deadlines.


	% \subsection{What happened during the scheduling of each scheduler}
	% \textbf{CFS.} The Completely Fair Scheduler (CFS) is not designed for real-time scheduling. Its primary goal is to allocate CPU time fairly among all tasks. When overall CPU utilization is low, CFS achieves a good balance in system performance. However, as SLAM's execution time increases and overall CPU utilization rises, CFS does not adapt to the tasks' differing importance or resource requirements. Instead, it continues to distribute resources evenly across tasks. Consequently, SLAM frequently misses its deadlines, leading to a gradual degradation of overall system performance.

	% \textbf{DM\_FAST, DM\_SLOW.} Deadline-Monotonic (DM) scheduling is a real-time algorithm that prioritizes tasks with shorter relative deadlines. In this case, TSP is assigned a higher priority than SLAM because of its shorter period.

    % In DM\_FAST, TSP is allocated the shortest possible running time limit, increasing the likelihood that all tasks meet their deadlines. However, this comes at the cost of TSP's performance metric, which is the lowest among the tested configurations.

    % In contrast, in DM\_SLOW, TSP is allocated a longer running time limit, improving its performance metric but making the overall system harder to schedule. While this initially leads to better performance for DM\_SLOW compared to DM\_FAST, due to SLAM's shorter initial running time, it becomes less effective as SLAM's execution time increases.

    % Both DM\_FAST and DM\_SLOW fail to adapt to the changes in the environments or the computation tasks, and therefore cannot exploit the full computation capabilities provided by the computing system.

	% \textbf{Optimizer-BF, incremental.}
    % The proposed optimizer addresses the limitations of the baseline schedulers by dynamically adjusting task priorities and running time limits in response to changes in resource requirements.
    % At the start of the experiments, when SLAM has a shorter execution time, the optimizer assigns higher priority and a longer running time limit to TSP, maximizing TSP's performance metric.
    % As SLAM's execution time increases, the optimizer reassigns higher priority to SLAM, ensuring it meets its deadlines, given its greater importance (as indicated by SLAM's weight parameter and safety threshold). 
	% Concurrently, TSP's running time limit is reduced based on system resource availability, optimizing overall system performance.
    % This adaptive approach not only improves overall performance compared to baseline methods but also ensures greater stability and robustness under dynamic conditions.

    % Readers may be curious whether the same effect can be achieved by always assigning high-criticality tasks higher priorities in the experiments.
	% Unfortunately, this simple method fails to consider schedulability during priority assignments, and may provide poor (computational) safety performance in some overloaded systems, especially if some low-criticality tasks have shorter running periods than some high-criticality tasks.



	% --- superseded by "Performance Comparison of Different Schedulers" (Beat 2) ---
	% \subsection{Proposed Schedulers' Performance Comparison}
	%
    % A noteworthy observation is that the brute-force optimizer sometimes performs worse than the incremental optimizer in terms of overall system performance. While the brute-force scheduler can find optimal solutions for priority assignments and task time limits, it has a significant drawback: long execution time. This drawback negatively impacts the system in two ways:
    % (1) The scheduler's long execution time increases its own resource demands and may cause competition in some computation resources such as memory and I/O. This may influence the primary tasks in high-utilization systems, which degrades performance.
    % (2) The delay in updating task priorities and configurations limits the system's ability to respond quickly to changes in task demands or environmental conditions.
    %
    % The incremental optimizer, on the other hand, delivers better and more robust performance because it dynamically adjusts task scheduling decisions in real-time, maintaining adaptability across a wide range of scenarios. This dynamic capability enables the incremental optimizer to outperform baseline schedulers, which are designed for static environments and are only effective under limited conditions.
	%
% As shown in Fig.~\ref{figs_exp_results_all_cps}, standard schedulers like CFS and DM fail to achieve optimal or stable performance in dynamic environments without significant modifications or system-level enhancements. Such static scheduling algorithms typically only work the best under specific situations.
% The proposed optimization-based scheduler, by contrast, demonstrates both improved performance and greater robustness, offering a practical solution for achieving near-optimal system performance in dynamic conditions.


	% \todo show DM/CFS overhead?

    % Section~\ref{section_complexity_analysis} shows that the incremental optimization algorithm has linear runtime complexity relative to the number of tasks, comparable to standard scheduling algorithms like DM and CFS. This ensures excellent scalability.

% While our experiments involve a limited number of tasks, we expect greater performance improvements between the incremental and the brute-force schedulers in systems with more tasks because:
% (1) Higher scheduling complexity: Larger systems offer more opportunities for optimization.
% (2) Low computational cost: The linear complexity ensures that scaling up remains cost-effective compared to baseline methods.
% This makes the incremental algorithm highly suitable for complex scheduling scenarios.


	\subsection{Performance Comparison of Different Schedulers}
    CFS fixes task priorities and allocates CPU time fairly without regard to
    importance or deadlines.
    When utilization is low this yields a reasonable balance, but as SLAM's
    execution time rises, CFS made SLAM
    frequently misses its deadlines and overall performance degrades.

    DM\_FAST and DM\_SLOW fix the QoS budget instead, using Deadline-Monotonic
    priorities that rank TSP above SLAM by period.
    DM\_FAST sets TSP the shortest budget, protecting schedulability but giving
    TSP the lowest performance among the tested configurations.
    DM\_SLOW sets TSP a longer budget, improving TSP's performance early on but
    becoming unschedulable as SLAM's execution time grows.
    Neither responds to the environment, so neither exploits the system's full
    computation capability.

    The proposed optimizer, in both its brute-force and incremental forms,
    adapts in both dimensions.
    It assigns higher priority and a longer QoS budget to TSP early in a run,
    when SLAM is cheap, maximizing TSP's performance; as SLAM's execution time
    rises it reassigns higher priority to SLAM and trims TSP's budget, keeping
    SLAM schedulable given its greater weight and safety threshold.
    This adaptive behavior yields higher and more stable SP than the baselines,
    as Fig.~\ref{figs_exp_results_all_cps} shows.

    % A natural question is whether the same effect comes from always assigning
    % the safety-critical tasks higher priorities.
    % It does not: such a rule ignores schedulability during priority assignment
    % and can give poor safety performance in overloaded systems, especially when
    % some low-criticality tasks have shorter periods than high-criticality ones.

    The optimizers' own cost also impacts performance.
    The brute-force optimizer's long execution time hurts it in two ways: its
    own resource demand competes with the primary tasks under high utilization,
    and the delay in updating configurations limits the system's responsiveness
    to demand changes.
    The incremental optimizer avoids both, running roughly an order of magnitude
    faster while holding linear complexity in the number of tasks, so it usually
    matches or exceeds the brute-force optimizer despite the latter enumerating
    the full joint priority-and-QoS space; we quantify the gap to the
    brute-force optimum rigorously in Section~\ref{section: simulation}
    (Table~\ref{tab_sim_vs_optimal}).


    % --- superseded by "Safety-Performance Guarantee" (Beat 3) ---
    % \subsection{Can we provide a safety-performance guarantee?}
    % /todo{UPDATE this section with updated real-world results}
    % % TODO: discuss the dependency of online safety based on ET prediction accuracy.
    % % so if ET chagnes very fast, then online ET prediction will not be accurate. So by
    %
    % In real-time systems, providing guarantees typically involves ensuring that a target metric—such as schedulability—is met under worst-case scenarios. These guarantees rely on offline analysis of worst-case conditions, as they serve as the foundation for establishing system performance bounds. Without such assumptions, guarantees become infeasible.
    %
    % This paper can provide safety-performance guarantee if the execution time distribution in the online environments can be upper-bounded and the estimation is known in advance.
    % Although this assumption is in some conflict with our original proposal about working in an unknown environment, there are no way to provide any performance guarantee if such worst-case information cannot be obtained offline.
    % However, the proposed optimization framework will still be useful in soft real-time systems where strict performance guarantee is not required.
    %
    % While existing scheduling methods like Completely Fair Scheduling (CFS) or Deadline-Monotonic (DM) can also provide guarantees within this framework, the proposed optimization framework achieves significantly higher SP metric values. This improvement underscores one of the key contributions of this work: a new, and stronger safety-performance guarantee that surpasses traditional methods in both safety and performance dimensions.
    %
    % Furthermore, the framework offers flexibility by allowing additional constraints to be incorporated for individual tasks. This enables the creation of fine-grained guarantees tailored to specific safety or performance requirements, making the approach adaptable to diverse real-time system needs.

    \subsection{Safety-Performance Guarantee}
    \label{section: real_sp_guarantee}
    The framework provides a formal schedulability guarantee for the
    safety-critical tasks $\boldsymbol{\tau}^{safe}$.
    By Theorem~\ref{obs:safety_fallback}, provided the online execution-time
    prediction upper-bounds the true distribution, every configuration the
    scheduler ships satisfies $Pr(r_i > D_i) \leq \Theta_i$ for each
    $\tau_i \in \boldsymbol{\tau}^{safe}$ at every invocation.

    The guarantee rests on two mechanisms in
    Section~\ref{section_safety_fallback}: the in-search feasibility gate,
    which rejects any candidate that would violate an important task's
    deadline-miss tolerance during optimization, and the certified safe
    fallback, an offline-verified DM-Fast configuration deployed whenever
    the inherited incumbent becomes unschedulable or an execution-time jump
    invalidates it.

    In our real-world runs the safety-critical task is MPC
    (Table~\ref{table_exp_config}).
    MPC met its deadline-miss tolerance ($\Theta = 0.99$) throughout all
    runs, consistent with Section~\ref{section_safety_fallback}, where
    monitoring abrupt execution-time changes suffices to keep
    safety-critical tasks schedulable without a hard worst-case
    execution-time bound online.

    Unlike fixed baselines such as CFS or DM, which can also be schedulable
    but sacrifice performance, the optimizer preserves this important-task
    guarantee while improving the SP metric---a stronger safety-performance
    guarantee than static scheduling provides.



\subsection{Reducing optimization overhead by changing its running frequency}
The scheduler runs once every $10$~s re-optimization interval rather than every
period, so its per-invocation cost is amortized over that interval.
In our real-world runs the incremental optimizer takes about $0.2$~s per
invocation, roughly $2\%$ of the interval
(Fig.~\ref{fig_et_all}, \ref{figs_exp_results_all_cps}).

This is a prototype implementation, not tuned for raw speed, and the measured
cost is not the focus of this work.
What matters for scalability is the algorithm's complexity: the incremental
optimizer runs in time linear in the number of tasks
(Section~\ref{section_complexity_analysis}), so the cost grows proportionally
rather than explosively as the task set expands, and lowering the running
frequency further reduces overhead whenever the environment changes slowly.

